Für $f\in F$ Zeigen Sie, dass die folgenden Identitäten gelten:
\begin{teilaufgaben}
\item $n\log(f) = \exp(f^n)$
\item $\exp(nf) = \exp(f)^n$
\item $\log(\exp(f)) = f$
\item $\exp(\log f)) = f$
\end{teilaufgaben}

\begin{hinweis}
Man betrachte in b) und d) die logarithmische Ableitung von Aufgabe
\ref{702}.
\end{hinweis}

\begin{loesung}
\begin{teilaufgaben}
\item
Die Ableitungen beider seiten sind
\begin{align*}
D(n\log(f) &= n\, Df/f
D\log(f^n) &= D(f^n)/f^n = nf^{n-1}\, Df/f^n = n\,Df/f
\end{align*}
Da die Differnz dieser Gleichungen
\[
D(n\log(f)-\log(f^n)) = 0
\]
ist, können sich die beiden Funktionen nur um eine Konstante unterscheiden.

\item
Anwendung des Operators $D$ auf beide Seiten liefert
\begin{align*}
D\exp(nf)  &= \exp(nf) n Df \\
D\exp(f)^n &= n \exp(f)^{n-1} D\exp(f) = n \exp(f)^n Df
\end{align*}
Die Funktionen $g_1=\exp(nf)$ und $g_2=\exp(f)^n$ erfüllen also beide
die Differentialgleichung
\[
Dg = g\cdot n\,Df,
\qquad\Rightarrow\qquad
\frac{Dg}{g} = n\,Df.
\]
Somit ist
\begin{align*}
\frac{D(g_1/g_2)}{g_1/g_2}
&=
\frac{g_2}{g_1}
\frac{Dg_1\cdot g_2-g_1\cdot Dg_2}{g_2^2}
=
\frac{Dg_1\cdot g_2 -g_1\cdot Dg_2}{g_1g_2}
=
\frac{Dg_1}{g_1} - \frac{Dg_2}{g_2}
=
0,
\end{align*}
die beiden Funktionen unterscheiden sich also nur um einen konstanten
Faktor.

\item
Die Ableitung der linken Seite ist
\begin{align*}
D\log(\exp(f)) &= \frac{D\exp(f)}{\exp(f)} = \frac{\exp(f) Df}{\exp(f)}=Df.
\end{align*}
Somit ist 
\[
D\bigl(\log(\exp(f)) - f\bigr) = 0,
\]
die beiden Seiten unterscheiden sich also nur um eine Konstante.

\item
Wir berechnen die logarithmische Ableitung der beiden Seiten der Gleichung:
\begin{align*}
D\exp(\log(f))
&=
\exp(\log(f)) D\log(f)
=
\exp(\log(f)) \frac{Df}{f}
\\
\Rightarrow\quad
\frac{D\exp(\log(f))}{\exp(\log(f))}
&=
\frac{Df}{f},
\end{align*}
die beiden Seiten der Gleichung können sich also nur um einen konstanten
Faktor unterscheiden.
\qedhere
\end{teilaufgaben}
\end{loesung}

